\subsection{Kỹ năng quản lý}

\subsubsection{Kỹ năng quản lý source code với Git/GitHub}
Hệ thống quản lý phiên bản phân tán giúp kiểm soát mã nguồn dự án chặt chẽ và an toàn:
\begin{itemize}
    \item \textbf{Chiến lược tạo nhánh (Branching):} Mỗi thành viên phát triển tính năng trên một nhánh riêng biệt gắn liền với mã số sinh viên hoặc tên nhiệm vụ, bảo vệ sự ổn định của nhánh \texttt{master}.
    \item \textbf{Quy trình Pull Request và Code Review:} Thực hiện mở PR để đề xuất tích hợp code vào nhánh chính; các thành viên khác rà soát chất lượng mã nguồn trước khi phê duyệt merge.
    \item \textbf{Kiểm soát lịch sử và giải quyết conflict:} Ghi chép commit message rõ ràng, theo dõi lịch sử chỉnh sửa và giải quyết thành thạo các xung đột mã nguồn trực tiếp trên giao diện GitHub.
\end{itemize}

\subsubsection{Kỹ năng quản lý task với Trello}
Trello hỗ trợ điều phối công việc trực quan theo mô hình Kanban:
\begin{itemize}
    \item \textbf{Trực quan hóa quy trình làm việc:} Phân chia các giai đoạn rõ ràng trên bảng: \textit{To Do} (Cần làm), \textit{In Progress} (Đang làm), \textit{Review/Testing} (Kiểm thử), và \textit{Done} (Hoàn thành).
    \item \textbf{Quản lý chi tiết từng đầu việc:} Mỗi tính năng hay lỗi phát sinh được tạo thành một Card riêng có gắn người thực hiện, checklist chi tiết và thời hạn hoàn thành (Due Date).
    \item \textbf{Kiểm soát tiến độ toàn dự án:} Giúp nhóm kịp thời phát hiện các công việc bị nghẽn để kịp thời phân bổ lại nguồn lực, đảm bảo dự án về đích đúng hẹn.
\end{itemize}